% =====================================================================
% CAP. 11 - NIVEL 3
% =====================================================================
\blocoaprofundamento

\begin{questao}[3]{}
Um capacitor de $\SI{10}{\micro\farad}$, inicialmente descarregado, é ligado a
$\SI{12}{\volt}$ por $R_1=\SI{2}{\kilo\ohm}$ durante
$\SI{50}{\milli\second}$. Em seguida, a fonte é removida e ele passa a
descarregar por $R_2=\SI{5}{\kilo\ohm}$. Determine sua tensão
$\SI{40}{\milli\second}$ após a segunda comutação.
\resposta{$v_C\approx\SI{4.95}{\volt}$}
\resolucao{Na primeira etapa, $\tau_1=\SI{20}{\milli\second}$ e
$v_C(50\,\mathrm{ms})=12(1-e^{-2{,}5})\approx\SI{11.015}{\volt}$. Esse valor é a
condição inicial da segunda etapa. Agora $\tau_2=\SI{50}{\milli\second}$ e
$v_C=11{,}015e^{-40/50}\approx\SI{4.95}{\volt}$.}
\end{questao}

\begin{questao}[3]{}
Um indutor de $\SI{100}{\milli\henry}$ é energizado por uma fonte de
$\SI{20}{\volt}$ através de $R_1=\SI{10}{\ohm}$ durante
$\SI{15}{\milli\second}$. Depois, a fonte é retirada e o indutor descarrega por
$R_2=\SI{25}{\ohm}$. Qual é a corrente após mais $\SI{8}{\milli\second}$?
\resposta{$i_L\approx\SI{0.210}{\ampere}$}
\resolucao{Na energização, $I_f=20/10=\SI{2}{\ampere}$ e
$\tau_1=0{,}1/10=\SI{10}{\milli\second}$. Assim
$i_L(15\,\mathrm{ms})=2(1-e^{-1{,}5})\approx\SI{1.554}{\ampere}$. Na descarga,
$\tau_2=0{,}1/25=\SI{4}{\milli\second}$, portanto após $2\tau_2$:
$i_L=1{,}554e^{-2}\approx\SI{0.210}{\ampere}$.}
\end{questao}

\begin{questao}[3]{}
Após a supressão das fontes independentes, a resistência vista por um
capacitor de $C=\SI{20}{\micro\farad}$ é composta por
$R_3=\SI{1}{\kilo\ohm}$ em série com o paralelo de
$R_1=\SI{6}{\kilo\ohm}$ e $R_2=\SI{3}{\kilo\ohm}$. Determine a
constante de tempo.

\circuitoq{%
\begin{circuitikz}[scale=1,transform shape,line width=0.8pt]
\draw (0,1.5) to[C,l={$C$},o-] (2,1.5) to[R,l={$R_3$}] (4,1.5) coordinate(a);
\draw (a) -- (4,2.6) to[R,l={$R_1$}] (6,2.6) -- (6,0.4);
\draw (a) -- (4,0.4) to[R,l_={$R_2$}] (6,0.4) -- (6,-0.2) node[ground]{};
\end{circuitikz}}

\resposta{$R_{\mathrm{eq}}=\SI{3}{\kilo\ohm}$; $\tau=\SI{60}{\milli\second}$}
\resolucao{$R_1\parallel R_2=(6\cdot3)/(6+3)=\SI{2}{\kilo\ohm}$. Em série com
$R_3$, $R_{\mathrm{eq}}=\SI{3}{\kilo\ohm}$. Logo
$\tau=R_{\mathrm{eq}}C=\pot{3}{3}\pot{20}{-6}=\SI{60}{\milli\second}$.}
\end{questao}

\begin{questao}[3]{}
Após uma comutação, $v_C(0^+)=-\SI{4}{\volt}$, $v_C(\infty)=\SI{8}{\volt}$,
$R_{\mathrm{eq}}=\SI{2.5}{\kilo\ohm}$ e $C=\SI{40}{\micro\farad}$. Em que
instante a tensão cruza zero?
\resposta{$t\approx\SI{40.5}{\milli\second}$}
\resolucao{$\tau=\SI{0.1}{\second}$ e
$v_C=8+(-4-8)e^{-t/0{,}1}=8-12e^{-t/0{,}1}$. No cruzamento,
$e^{-t/0{,}1}=2/3$. Portanto $t=0{,}1\ln(3/2)\approx\SI{40.5}{\milli\second}$.}
\end{questao}

\begin{questao}[3]{}
Um indutor possui $i_L(0^+)=-\SI{1}{\ampere}$, $i_L(\infty)=\SI{3}{\ampere}$ e
$\tau=\SI{25}{\milli\second}$. Determine quando a corrente muda de sinal.
\resposta{$t\approx\SI{7.19}{\milli\second}$}
\resolucao{$i_L=3+(-1-3)e^{-t/\tau}=3-4e^{-t/\tau}$. Para $i_L=0$,
$e^{-t/\tau}=3/4$. Logo $t=\tau\ln(4/3)\approx\SI{7.19}{\milli\second}$.}
\end{questao}

\begin{questao}[3]{}
Um capacitor com $v_C(0)=\SI{2}{\volt}$ é ligado a uma fonte de
$\SI{10}{\volt}$ através de $R=\SI{2}{\kilo\ohm}$ e
$C=\SI{50}{\micro\farad}$. Determine a corrente inicial e a corrente após
$\SI{150}{\milli\second}$.
\resposta{$i(0^+)=\SI{4}{\milli\ampere}$; $i(150\,\mathrm{ms})\approx\SI{0.893}{\milli\ampere}$}
\resolucao{$\tau=RC=\SI{0.1}{\second}$. A corrente é
$i(t)=(10-2)R^{-1}e^{-t/\tau}=\SI{4}{\milli\ampere}e^{-t/\tau}$. Em
$150\,\mathrm{ms}=1{,}5\tau$, $i\approx4e^{-1{,}5}=\SI{0.893}{\milli\ampere}$.}
\end{questao}

\begin{questao}[3]{}
Um capacitor de $\SI{100}{\micro\farad}$ inicialmente a $\SI{20}{\volt}$ é
descarregado por $R=\SI{1}{\kilo\ohm}$. Mostre, por integração da potência no
resistor, que toda a energia inicialmente armazenada é dissipada e calcule-a.
\resposta{$E=\SI{20}{\milli\joule}$}
\resolucao{$v_C=V_0e^{-t/RC}$ e $i=(V_0/R)e^{-t/RC}$. Logo
$p_R=Ri^2=(V_0^2/R)e^{-2t/RC}$. Integrando de zero a infinito:
\[
E_R=\frac{V_0^2}{R}\int_0^\infty e^{-2t/RC}\,\mathrm{d}t
=\frac{V_0^2}{R}\frac{RC}{2}=\frac12CV_0^2.
\]
Numericamente, $E_R=\tfrac12\pot{100}{-6}\cdot20^2=\SI{20}{\milli\joule}$.}
\end{questao}

\begin{questao}[3]{}
Projete um atraso RC para que uma tensão de $\SI{5}{\volt}$ alcance
$\SI{3.3}{\volt}$ em $\SI{0.5}{\second}$. Use $R=\SI{100}{\kilo\ohm}$ e
considere o capacitor inicialmente descarregado. Determine $C$ e indique um
valor comercial próximo.
\resposta{$C\approx\SI{4.63}{\micro\farad}$; valor comercial: \SI{4.7}{\micro\farad}}
\resolucao{$3{,}3=5(1-e^{-t/RC})$, então
$C=-t/[R\ln(1-3{,}3/5)]$. Substituindo,
$C\approx\SI{4.63}{\micro\farad}$. Um valor comercial próximo é
\SI{4.7}{\micro\farad}.}
\end{questao}

\begin{questao}[3]{}
Deseja-se que a corrente de um circuito RL alcance \SI{95}{\percent} de seu valor
final de $\SI{1}{\ampere}$ em $\SI{20}{\milli\second}$. Se
$L=\SI{50}{\milli\henry}$, determine a resistência e a tensão da fonte.
\resposta{$R\approx\SI{7.49}{\ohm}$; $V\approx\SI{7.49}{\volt}$}
\resolucao{$0{,}95=1-e^{-t/\tau}$ implica
$\tau=-t/\ln0{,}05\approx\SI{6.68}{\milli\second}$. Como $\tau=L/R$,
$R=0{,}05/0{,}00668\approx\SI{7.49}{\ohm}$. Para corrente final de \SI{1}{\ampere},
$V=RI_f\approx\SI{7.49}{\volt}$.}
\end{questao}

\begin{questao}[3]{}
Um RLC em série tem $R=\SI{20}{\ohm}$, $L=\SI{0.1}{\henry}$ e
$C=\SI{100}{\micro\farad}$. Determine os polos naturais.
\resposta{$s_{1,2}=-100\pm j300\ \si{\per\second}$}
\resolucao{$\alpha=R/(2L)=100$ e $\omega_0\approx316{,}2$. Assim
$\omega_d=\sqrt{316{,}2^2-100^2}=300$. Para o caso subamortecido,
$s_{1,2}=-\alpha\pm j\omega_d=-100\pm j300\ \si{\per\second}$.}
\end{questao}

\begin{questao}[3]{}
Uma resposta criticamente amortecida é
$v(t)=(A+Bt)e^{-200t}$. Se $v(0)=\SI{10}{\volt}$ e
$\dot v(0)=-\SI{500}{\volt\per\second}$, determine $A$ e $B$.
\resposta{$A=\SI{10}{\volt}$; $B=\SI{1500}{\volt\per\second}$}
\resolucao{De $v(0)=A$, segue $A=10$. Derivando,
$\dot v(0)=B-200A=-500$. Logo $B=-500+2000=\SI{1500}{\volt\per\second}$.}
\end{questao}

\begin{questao}[3]{}
Para um RLC em série com $R=\SI{100}{\ohm}$, $L=\SI{0.1}{\henry}$ e
$C=\SI{100}{\micro\farad}$, determine os dois polos reais.
\resposta{$s_1\approx-\SI{112.7}{\per\second}$; $s_2\approx-\SI{887.3}{\per\second}$}
\resolucao{A equação característica é
$s^2+(R/L)s+1/(LC)=s^2+1000s+100000=0$. Portanto
\[
s=\frac{-1000\pm\sqrt{1000^2-4\cdot100000}}{2}
=\frac{-1000\pm\sqrt{600000}}{2},
\]
resultando aproximadamente em $-112{,}7$ e $-887{,}3\ \si{\per\second}$.}
\end{questao}

\begin{questao}[3]{}
Determine a resistência crítica de um RLC \textbf{paralelo} com
$L=\SI{50}{\milli\henry}$ e $C=\SI{200}{\micro\farad}$.
\resposta{$R_{\mathrm{crit}}\approx\SI{7.91}{\ohm}$}
\resolucao{No circuito RLC em paralelo, a condição crítica é
$1/(2RC)=1/\sqrt{LC}$. Logo
$R_{\mathrm{crit}}=\tfrac12\sqrt{L/C}
=\tfrac12\sqrt{0{,}05/\pot{200}{-6}}\approx\SI{7.91}{\ohm}$.}
\end{questao}

\begin{questao}[3]{}
Para um sistema subamortecido de segunda ordem com $\zeta=0{,}5$, estime o
sobressinal percentual de pico usando
$M_p=e^{-\pi\zeta/\sqrt{1-\zeta^2}}$.
\resposta{$M_p\approx\SI{16.3}{\percent}$}
\resolucao{Substituindo $\zeta=0{,}5$:
$M_p=e^{-\pi(0{,}5)/\sqrt{0{,}75}}\approx0{,}163$. Portanto o sobressinal é
aproximadamente \SI{16.3}{\percent}.}
\end{questao}

\begin{questao}[3]{}
Um RLC em série possui $R=\SI{20}{\ohm}$, $L=\SI{100}{\milli\henry}$ e
$C=\SI{100}{\micro\farad}$. Calcule o fator de amortecimento $\zeta$.
\resposta{$\zeta\approx0{,}316$}
\resolucao{$\alpha=R/(2L)=100$ e $\omega_0\approx316{,}2$. Assim
$\zeta=\alpha/\omega_0\approx100/316{,}2\approx0{,}316$.}
\end{questao}

\begin{questao}[3]{}
Classifique os quatro casos da tabela.

\begin{table}[H]
\centering\small\sffamily
\begin{tabular}{@{}c c c c c@{}}
\toprule
\textbf{Caso} & \textbf{Topologia} & $R$ & $L$ & $C$\\
\midrule
A & série & \SI{20}{\ohm} & \SI{0.1}{\henry} & \SI{100}{\micro\farad}\\
B & série & \SI{63.25}{\ohm} & \SI{0.1}{\henry} & \SI{100}{\micro\farad}\\
C & série & \SI{100}{\ohm} & \SI{0.1}{\henry} & \SI{100}{\micro\farad}\\
D & paralelo & \SI{20}{\ohm} & \SI{50}{\milli\henry} & \SI{200}{\micro\farad}\\
\bottomrule
\end{tabular}
\end{table}
\resposta{A subamortecido; B aproximadamente crítico; C superamortecido; D subamortecido}
\resolucao{A, B e C possuem $\omega_0\approx316{,}2$. No circuito RLC em série,
$\alpha=R/(0{,}2)$: A dá $100<316{,}2$; B dá aproximadamente $316{,}25$; C dá
$500>316{,}2$. No caso D, $\alpha=1/(2RC)=125$ e
$\omega_0=1/\sqrt{0{,}05\cdot\pot{200}{-6}}\approx316{,}2$, logo também é
subamortecido.}
\end{questao}

\begin{questao}[3]{}
Em certo instante, o capacitor de $C=\SI{100}{\micro\farad}$ apresenta
tensão de $\SI{10}{\volt}$, enquanto o indutor de
$L=\SI{50}{\milli\henry}$ conduz corrente de $\SI{0.2}{\ampere}$. Qual é a energia total armazenada naquele instante?
\resposta{$E=\SI{6}{\milli\joule}$}
\resolucao{$E_C=\tfrac12CV^2=\SI{5}{\milli\joule}$ e
$E_L=\tfrac12LI^2=\tfrac12(0{,}05)(0{,}2)^2=\SI{1}{\milli\joule}$. A soma é
\SI{6}{\milli\joule}.}
\end{questao}

\begin{questao}[3]{}
Antes de uma comutação, um capacitor possui $v_C=\SI{6}{\volt}$ e um indutor
possui $i_L=\SI{0.3}{\ampere}$. Logo após a comutação, ambos passam a integrar
uma nova rede resistiva. Quais condições iniciais devem ser usadas na nova rede?
\resposta{$v_C(0^+)=\SI{6}{\volt}$ e $i_L(0^+)=\SI{0.3}{\ampere}$}
\resolucao{As variáveis de estado associadas aos elementos ideais armazenadores
são contínuas: a tensão do capacitor e a corrente do indutor mantêm seus valores
no instante da comutação. A nova topologia altera suas derivadas e seus valores
finais, não os valores instantâneos em $0^+$.}
\end{questao}

\begin{questao}[3]{}
Para um RLC em série alimentado por um degrau de $\SI{10}{\volt}$, use
$R=\SI{20}{\ohm}$, $L=\SI{0.1}{\henry}$ e $C=\SI{100}{\micro\farad}$. Escreva
a equação diferencial para $v_C(t)$ após a comutação.
\resposta{$\ddot v_C+200\dot v_C+100000v_C=1000000$}
\resolucao{Pela LKT,
$LC\ddot v_C+RC\dot v_C+v_C=V_s$. Como $LC=\pot{1}{-5}$ e
$RC=\pot{2}{-3}$, dividindo por $LC$ resulta
\[
\ddot v_C+200\dot v_C+100000v_C=10\cdot100000=1000000.
\]}
\end{questao}

\begin{questao}[3]{}
No circuito da questão anterior, admita energia inicial nula. Determine
$v_C(0^+)$, $\dot v_C(0^+)$ e $\ddot v_C(0^+)$.
\resposta{$v_C(0^+)=0$; $\dot v_C(0^+)=0$; $\ddot v_C(0^+)=\SI{1e6}{\volt\per\second\squared}$}
\resolucao{Energia inicial nula implica $v_C(0^+)=0$ e $i_L(0^+)=0$. Como
$i_C=C\dot v_C$, segue $\dot v_C(0^+)=0$. Substituindo esses valores na equação
$\ddot v_C+200\dot v_C+100000v_C=1000000$, obtém-se
$\ddot v_C(0^+)=\SI{1e6}{\volt\per\second\squared}$.}
\end{questao}

\gabarito[Nível 3]
